\section{Innovation Is a Game}\label{sec:game}
% Reader question: Can invention be played like Go?
% Core points advanced: 1 (primary: innovation is a game), 2 (novelty is free;
% judgment is the inventor), 3 (search = 1)
% Exhibits: Table 4 (tab:go-vs-innovation), Figure 3 (fig:game-loop)

Can invention be played like Go? It can. The problem is the board, an
algorithm is the move, a verifier is the referee and the best known algorithm
is the opponent. Because the problem sits on the board rather than in the
rules, one set of rules covers many kinds of problem.

% ----------------------------------------------------------------------------
\subsection{Why games came first}\label{sub:game-games}

Games became AI's laboratory because they supply four things at once: fixed
rules, an impartial referee, a clear winner and cheap replay. Chess was chosen
because it is ``sharply defined both in allowed operations (the moves) and in
the ultimate goal (checkmate)'' \citep{shannon1950chess}, and checkers showed
that a program could improve by playing itself \citep{samuel1959checkers}.
Together the four give a learner unlimited experience, scored honestly. People
scale judgment the same way, through prizes, leaderboards and the common task
framework's ``scoring referee'' (\citealp{donoho2017fifty}; \cref{sec:road}).
Invention has the first three properties once the problem and the solution are
written down; cheap replay is a question of cost (\cref{sec:price}).

% ----------------------------------------------------------------------------
\subsection{The rules}\label{sub:game-rules}

\Cref{tab:go-vs-innovation} sets the game beside Go. Two rows differ in kind:
an illegal move is played and lost, and the board holds the incumbent as well as
the problem, because a move can only win against something.

\begin{table}[htb]
  \caption{\textbf{Invention has every part of a board game, except cheap
  moves.} Go and the innovation game, element by element. Go's game length is
  from \citet{silver2016mastering}.}
  \label{tab:go-vs-innovation}
  \centering
  \small
  \begin{tabular}{@{}p{0.16\linewidth}p{0.33\linewidth}p{0.43\linewidth}@{}}
    \toprule
                   & Go                              & The innovation game \\
    \midrule
    Board          & Stones on a $19\times19$ grid   & The encoded problem plus the incumbent \\
    Move           & Place one stone                 & Write one complete program, or one edit \\
    Legal move     & An empty point; no suicide      & Parses, type-checks, runs within limits \\
    Illegal move   & Not allowed                     & Stillborn: an automatic loss \\
    Referee        & The rules of Go                 & The verifier, on hidden folds \\
    Win            & More territory                  & Valid, and beats the incumbent \\
    Opponent       & A past self                     & The incumbent, often a past self \\
    Rating         & Elo                             & Win rate and Elo across problems \\
    Moves per game & About 150, each cheap           & One, costly to judge \\
    \bottomrule
  \end{tabular}
\end{table}

Formally, the game is a tuple of seven parts,
\[
  \Game = (\Prob, \encP, \Sol, \encS, \referee, \reward, \opp).
\]
Here $\encP$ and $\encS$ embed a problem $p \in \Prob$ and a program
$w \in \Sol$ as $z_p$ and $z_w$. The referee $\referee$ reports a program's
validity and score, the reward $\reward$ orders moves, and the opponent $\opp$
supplies the incumbent $\inc{p}$. The tuple extends Rice's model of algorithm
selection \citep{rice1976algorithm}: the player writes the algorithm instead of
picking it, and an opponent decides what counts as winning.

A skeptic might ask whether a one-move game is a game at all. Formally it is a
contextual bandit, with the board as context, the program as action and the
verdict as reward \citep{lattimore2020bandit}. It earns the name through an
opponent that improves as the player does. The honest asymmetry is in the
moves: Go offers about 250 legal moves per turn over about 150 turns, each free
to check \citep{silver2016mastering}, while this game offers one move, costly to
judge. Taste, which raises the chance $q$ that one try wins (\cref{eq:taste}),
therefore decides it.

\begin{figure}[t]
  \centering
  \includegraphics[width=\linewidth]{figures/fig3_game_loop.pdf}
  \caption{\textbf{Training plays the whole game; game day plays one move.}
  Schematic. Top row (blue): game day. Dashed box: training only, where every
  verdict, losses included, becomes a lesson.}
  \label{fig:game-loop}
\end{figure}

\Cref{fig:game-loop} draws the loop. Game day plays only its top row, so
nothing can send a move back for another try; that is \searchone.

% ----------------------------------------------------------------------------
\subsection{The board and the move}\label{sub:game-board}

The board is the aim pillar of \cref{sec:anatomy}, made machine-readable: a
summary of the training data, the objective and metric, and the constraints.
The summary is DOPE's landmark encoding (\cref{sec:method}), in the
meta-learning tradition of describing a dataset by how cheap learners score on
it \citep{pfahringer2000landmarking}. A good board is \emph{complete}, since
problems that encode alike must be answered alike; \emph{canonical} when rows
and columns are reordered; \emph{leak-free}, blind to test labels and hidden
folds; and \emph{cheap}, at \TBD{seconds} per problem. Whoever writes the board
writes the aim. People still choose the problem, and the machine learns to
answer it.

The move is where invention gets its language: one program in a typed language
over learning primitives \TBD{primitive set and its size}. Types remove whole classes of stillborn
programs before they run \citep{montana1995stgp}. A canonical form makes two
spellings of one idea count once, as AutoML-Zero reused the score of any program
whose outputs matched one already scored \citep{real2020automlzero}. The
language bounds what can be invented: in the terms of \citet{boden1990creative},
it permits exploring its space to any depth but never transforming it. Tighter
languages waste fewer moves and looser ones say more (\cref{sec:price}).

% ----------------------------------------------------------------------------
\subsection{The referee and the opponent}\label{sub:game-referee}

The reward is an ordering first and a number second. Validity is a gate: any
stillborn program scores below any valid one. Past the gate, programs rank by
improvement over the incumbent on hidden folds, with ties within the referee's
noise going to the simpler program. Novelty never pays for losing: any novelty
bonus (\cref{sec:price}) must respect the order. For a stillborn program
$\tilde{w}$, an also-ran $w^-$ and a win $w^+$ on board $p$,
\[
  \reward(p, \tilde{w}) \;<\; \reward(p, w^-) \;<\; \reward(p, w^+).
\]

The opponent is a ratcheting incumbent. A program is promoted only on a clear
win, by \TBD{margin} across \TBD{number of} folds and then re-verified
(\cref{sec:method}), as AlphaGo Zero promoted a network only after it won more
than 55\% of evaluation games \citep{silver2017mastering}. The ratchet also
makes this self-play.
Subtracting the incumbent's score does not change which move is best; what
matters is what promotion changes. It raises the bar, turns yesterday's winners
into today's also-rans and so into hard negatives for the problem--program map
(\cref{sec:method}), and shifts which problems are worth playing.

% ----------------------------------------------------------------------------
\subsection{One game, many boards}\label{sub:game-boards}

Binary and multiclass classification, regression, ranking, forecasting and
relational prediction differ only in what the board says: the data summary, the
metric and the type of output. Beyond machine learning, a board needs only a referee, such as tests
for sorting \citep{mankowitz2023alphadev}, automatic evaluators for mathematical
constructions \citep{romeraparedes2024funsearch} and timers for accelerator
kernels \citep{novikov2025alphaevolve}. These are boards for the recipe, not
results of this paper (\cref{sec:meaning}).

Casting algorithm discovery as a game is not new, but earlier games wrote one
family of problems into their rules, as AlphaTensor did for matrix
multiplication and AlphaDev for sorting
\citep{fawzi2022alphatensor,mankowitz2023alphadev}. AlphaZero used one algorithm
for three games and trained a separate network for each
\citep{silver2018general}. Here the problem is an input on the board, so one
player can learn many boards.

\begin{keyidea}
Every discipline with a referee is a game no one has played yet.
\end{keyidea}

% ----------------------------------------------------------------------------
\subsection{One game played, and where the analogy breaks}\label{sub:game-played}

The running example is the sealed problem chosen by the paper's stated rule,
the median win (\cref{sec:results}): \TBD{dataset, task type, metric}. The
player encodes it in \TBD{time} and writes one program of \TBD{size}, which
passes its checks. Fit in \TBD{time}, the program scores \TBD{score} on hidden
folds against the incumbent's \TBD{score}: \TBD{a win or an also-ran}. Had it
crashed, the game would simply have been lost, with no repair and no second try.
The update is skipped because the problem is sealed: the verdict is reported,
never learned from.

The analogy breaks in four places. \textbf{The referee is noisy}: picking a
winner by a noisy validation score biases that score upward
\citep{cawley2010overfitting}, so promotion needs a margin. \textbf{Moves are expensive}, each a fit and a cross-validation
(\cref{sec:price}). \textbf{Judges can be gamed}, as METR documented for
frontier models in June 2025 \citep{metr2025rewardhacking}; \cref{sec:method}
lists our safeguards. And \textbf{no algorithm wins everywhere}: averaged
uniformly over all problems, no learner has lower error off the training set
than another, ours included \citep{wolpert1996nfl}. Real problems are structured, though, and a player that
reads the board can choose its bias board by board. \Cref{sec:results} tests
that bet against frontier-invented algorithms that run unchanged on every
problem, with head-to-head win rates beside Elo \citep{elo1978rating}, which
assumes strength is transitive \citep{balduzzi2018reevaluating}.

A game anyone can play is only useful if someone can afford to play it.
